\section{Introduction}
\label{sec:introduction}

Melt-pool geometry is a practical link between laser processing conditions and the thermal models used to describe laser powder bed fusion (LPBF). The fused boundary constrains where melting occurred, while solidification analysis requires temperature gradients and the cooling histories experienced by material. A model that reproduces a track boundary can therefore support geometric process assessment. Spatially and temporally resolved thermography also provides observations of pulsed heating and cooling during solidification~\cite{hooper2018}. Geometric calibration supplies a boundary constraint, while the temperature and cooling behavior associated with that boundary remains dependent on source and material assumptions and requires independent assessment.

Moving-source heat transfer provides a useful foundation for addressing this question at manageable computational cost. The progression from concentrated-source solutions to distributed heating resolves the singular near-source behavior of classical solutions, while nonlinear conduction formulations incorporate temperature-dependent properties and latent heat~\cite{vanelsen2007}. Rapid transient calculations extend this approach to scan-path effects and have been compared with solidification conditions and grain-structure trends in selected laser and electron-beam powder-bed cases~\cite{Plotkowski2017Rapid}. Their value lies in connecting process inputs to thermal conditions without resolving every melt-scale phenomenon. Fluid motion and deformation of the absorbing surface, however, redistribute energy through processes outside a conduction balance. Simplified models therefore rely on the heat-source representation and effective transport assumptions to reproduce the observed melt boundary~\cite{dynamic2024}.

Geometric calibration demonstrates both the usefulness and the coupling of these assumptions. In the AMMT IN625 benchmark, imposing the measured laser profile reduced the freedom to adjust the source shape. A scalar absorbed-power factor then reproduced length more closely than width and depth; introducing calibrated directional conductivity brought the dimensions into closer agreement across the tested conditions~\cite{kollmannsberger2019}. This result shows that a well-constrained incident beam does not remove uncertainty in how heat is redistributed within a simplified model. An alternative approach represents unresolved laser--melt interactions through a volumetric source whose deposition depth and absorption respond to the calculated pool. For bare IN625 tracks, calibration across several beam diameters supported subsequent application to a layer-scale benchmark~\cite{dynamic2024}. Together, these approaches establish that the calibrated geometry depends on the combination of source deposition and transport closure. Their fitted parameters consequently acquire meaning within that combination.

Temperature measurements provide a direct test of the thermal information supplied by such fits. For bare 316L tracks, several combinations of absorption and vaporization parameters in a fluid model gave comparable agreement with cross-sectional area and width-to-depth ratio, yet produced different surface-temperature profiles; two-color imaging distinguished among them~\cite{myers2023}. The authors of the AMMT conduction study explicitly stated that their calibrated model was not valid for predicting the temperature distribution within the melt pool~\cite{kollmannsberger2019}. Geometric observations also constrain different features of a field. In the AMMT measurements, imaging length uses radiance-profile features at the front and rear, whereas width and depth come from metallographic cross-sections and separate observation populations~\cite{lane2020}. Calibration to one of these quantities should thus be interpreted in terms of its measurement definition and its sensitivity to the modeled field.

The material description supplies another basis for that interpretation. Conductivity controls diffusive heat redistribution, while sensible heat capacity enters the enthalpy stored during heating and released during cooling. These inputs are often available over only part of the temperature range reached by a laser track. For the IN625 supplier tables used here, conductivity ends at 982~$^\circ$C and sensible specific heat at 1093~$^\circ$C, below the listed 1290--1350~$^\circ$C melting range. The bulletin labels conductivity as measured and specific heat as calculated typical alloy data~\cite{specialmetals625}. A continuation is therefore required before either table can describe the phase-change interval or the hotter liquid. Broader assessments of nickel-superalloy properties also show that liquid-property estimates depend on the available measurements and estimation procedure~\cite{mills2006}. A geometric fit cannot determine these constitutive functions independently of the source and transport model.

Continuing sensible properties must also be distinguished from specifying phase change. In an enthalpy formulation, latent heat contributes through the adopted phase-fraction relation as well as the total heat of fusion~\cite{vanelsen2007}. A linear fraction across a selected freezing interval is an established continuum approximation~\cite{dynamic2024}, but the interval itself depends on the material description. The supplier melting range supplies one set of endpoints; an IN625 model intended for rapid additive-manufacturing cooling instead terminated solidification at a lower eutectic temperature derived from Gulliver--Scheil calculations, motivated by limited diffusion in the solid~\cite{specialmetals625,dynamic2024}. These choices alter where latent storage is distributed along the temperature axis. They explain why a sensitivity to sensible heat capacity cannot be interpreted as an isolated latent-heat effect, and why comparisons of phase-region widths require a common phase law.

With that law fixed, a further distinction becomes available within the calculated temperature field. Total pool length combines the positions of the front and rear melting boundaries. At the rear, solidus and liquidus crossings describe both the location of the thermal tail and the distance over which material traverses the adopted phase interval. Moving both crossings downstream can lengthen the pool with little change in their separation. Changing their separation instead changes the spatial extent of that interval. Conductivity and storage influence the same coupled field, so a dimensional response alone does not identify which of these changes occurred. Resolving paired phase boundaries gives the geometric sensitivity a more specific thermal interpretation.

Converting that spatial structure to a cooling history requires the motion of material relative to the field. For a steadily translating temperature distribution, a stationary laboratory point traverses a rear spatial interval in a time equal to its length divided by scan speed. The constant-speed conversion is also used in the AMMT below-solidus cooling estimates~\cite{lane2020}. It follows that similar spatial phase spans can correspond to different passage times and interval-mean cooling rates when speed changes. The steady mapping has a defined physical scope: pulsed LPBF measurements show time histories affected by successive heating events and post-pulse latent release~\cite{hooper2018}, while fluid-tracer calculations show that melt motion changes the thermal paths followed by liquid parcels~\cite{Hou2024Dissolution}. Material passage through a translating conduction field therefore provides a useful thermal descriptor whose interpretation depends on the frame and transport assumptions.

The present study examines these distinctions for three AMMT bare-plate IN625 conditions using a nonlinear moving-frame enthalpy model. One effective source factor is fitted to condition B surface length; width, depth and conditions A and C are compared using the same factor. A controlled comparison then replaces linear continuation by last-value holding independently for conductivity and sensible heat capacity, while keeping the source, phase law and numerical configuration fixed. The analysis separates the resulting length changes into front and rear boundary motion, resolves rear solidus--liquidus displacement and span, and relates the spatial span to material passage time. This design identifies how the chosen property continuations change the thermal meaning of a length-calibrated field. It connects geometric discrepancy, constitutive response and the derived cooling history to distinct observables that can guide further thermal assessment.
